% Modernised reading — fonds Alexandre Grothendieck, shelfmark 54, pages 1-7.
%
% This is an INTERPRETATION, not a transcription. Everything the critical
% apparatus carried moves into footnotes. In any dispute of substance the
% transcription governs: whoever cites Grothendieck must cite batch-01.fr.tex.
%
% Pass: Opus 5 (claude-opus-5), 2026-08-30 — first pass, unchecked against the
% pages by a human. The folder was read whole; it holds one batch, pages 1-7.
%
% Notation fixed for the whole document, since the pages use three letters for
% two things: d = dim A throughout (page 2 writes n for the same integer,
% page 4 writes d); omega_A = H^0(A, Omega^1) and t_A = Lie A = omega_A^dual;
% A[p] for the manuscript's subscript-left {}_pA; ker F_A for its F^A, which it
% also writes gr(t_A) on page 4 and Gr(t_A) on page 6. The caron of page 2
% (t-check-B) is read as the linear dual t_B^dual = omega_B; nothing on the
% page defines it, and the reading says so in a footnote.
%
% Substantive departures from the pages, each footnoted where it occurs:
% — page 4 writes the graded duality as H^p(A, Omega^q_A) dual to
%   H^q(B, Omega^p_A). The second subscript must be B; with A it is false. The
%   exchange of p and q that the page does perform is right.
% — page 2's S_A and S_B are used without ever being defined. They are read as
%   the Hodge filtration sequences of A and of B, which is the only reading
%   under which the two claims made about them hold.
% — the proof of the corrected Proposition 4.1 of page 3 is not on these
%   leaves: the typed argument proves the nilpotent case, which is the case his
%   pen removes. The reading states the corrected statement and says so.
%
% What the folder announces and never establishes: the « structure commune »
% of page 4 that would define H*(A) and H*(B) at once and make their duality a
% formality — he writes that it « ne saurait offrir de difficultés » and does
% not write it. And the « conjugaison » of page 4 between the vector-extension
% picture and the Cartier picture is announced, worked at for two pages, and
% left at a question mark.
\documentclass[11pt,a4paper]{article}
% Le préambule commun à toutes les transcriptions du fonds Grothendieck.
%
% Il tient en une idée : **l'incertitude se compose, elle ne se résout pas.**
% Une transcription qui lisse un mot douteux a détruit la seule chose pour
% laquelle on la faisait — un lecteur doit pouvoir savoir, à chaque endroit,
% ce qui était sur le feuillet et ce qui a été ajouté. Les six macros qui
% suivent sont donc l'appareil critique, et non de la décoration.
%
% Les mêmes commandes sont reconnues par `scripts/render.mjs`, qui produit la
% vue de lecture du site. Une macro ajoutée ici doit l'être aussi là-bas,
% faute de quoi le rendu échouera bruyamment — ce qui est voulu : un rendu
% silencieusement faux est pire qu'un rendu absent.

\ProvidesPackage{grothendieck}[2026/08/08 Transcriptions du fonds Grothendieck]

\usepackage{amsmath,amssymb,amsthm}
\usepackage{mathrsfs}
\usepackage{stmaryrd}
% Les diagrammes commutatifs sont partout dans ce fonds : c'est la langue
% même de la géométrie algébrique de Grothendieck, et une transcription qui
% les rendrait en prose ne transcrirait rien.
\usepackage{tikz-cd}
% Français par défaut — c'est la langue des feuillets — mais l'anglais est
% chargé aussi : la traduction et le résumé sont des documents anglais, et la
% typographie française (espaces avant les deux-points) y serait une faute.
% Ces documents basculent par \selectlanguage{english} en tête de corps.
\usepackage[english,french]{babel}
\usepackage{fontspec}
\usepackage{geometry}
\geometry{a4paper,margin=25mm,marginparwidth=28mm}
\usepackage{marginnote}
\usepackage{xcolor}
% ulem, not soul. `soul`'s \st and \ul are built on a character-by-character
% scan that XeLaTeX's font machinery breaks: the document fails with an
% undefined control sequence rather than degrading. ulem does the same two jobs
% and is Unicode-engine safe. `normalem` stops it hijacking \emph, which the
% transcriptions use for Grothendieck's own emphasis.
\usepackage[normalem]{ulem}
\usepackage{hyperref}
\hypersetup{colorlinks=true,linkcolor=black,urlcolor=black,pdfborder={0 0 0}}

% --- Métadonnées du lot -----------------------------------------------------
% Reprises telles quelles par le script de rendu, qui les affiche en tête de
% la vue de lecture. Elles fixent ce que le fichier prétend couvrir : sans
% elles, une transcription flotte sans point d'attache dans un fonds de seize
% mille feuillets.

\newcommand{\folder}[1]{\gdef\grothFolder{#1}}
\newcommand{\batch}[1]{\gdef\grothBatch{#1}}
\newcommand{\leaves}[2]{\gdef\grothFirst{#1}\gdef\grothLast{#2}}
\newcommand{\foldertitle}[1]{\gdef\grothFoldertitle{#1}}
\gdef\grothFolder{?}\gdef\grothBatch{?}\gdef\grothFirst{?}\gdef\grothLast{?}\gdef\grothFoldertitle{}

% --- Le résumé --------------------------------------------------------------
%
% Il ouvre la lecture modernisée, sous le titre « Résumé », et il est composé
% en petit. Ce n'est pas un ornement : c'est le seul passage du document qui
% ne soit pas des mathématiques, et un lecteur doit pouvoir le reconnaître
% avant de l'avoir lu — puis le sauter s'il connaît déjà le sujet, ou s'y
% arrêter s'il ne le connaît pas. Le filet qui le suit marque la reprise du
% corps.

\newenvironment{resume}
  {\small\setlength{\parskip}{0.45em}}
  {\par\medskip\hrule\medskip}

% --- Mots-clés ---------------------------------------------------------------
%
% En anglais, en clôture du résumé de la lecture modernisée : le vocabulaire
% moderne sous lequel chercher ce que les feuillets construisent. Le manifeste
% du site (`npm run manifest`) les extrait pour étiqueter la cote entière —
% cette ligne est la source unique des tags, il n'y a pas de fichier à côté.
\newcommand{\keywords}[1]{\par\smallskip\noindent
  {\footnotesize\textsc{Keywords} --- \textit{#1}}\par}

% --- L'appareil critique ----------------------------------------------------

% Le numéro de feuillet, en marge. C'est l'ancre qui permet de régler un
% désaccord sur une formule en regardant *un* feuillet plutôt qu'une chemise
% de six cents. Le site s'en sert aussi pour tourner les pages du fac-similé
% quand on fait défiler la transcription.
\newcommand{\leaf}[1]{%
  \marginnote{\footnotesize\color{gray!70}\textsf{#1}}%
  \label{leaf:#1}\ignorespaces}

% Illisible : on ne devine pas. Un mot inventé qui se lit comme les autres est
% le pire résultat possible.
\newcommand{\ill}{\textcolor{red!60!black}{[\,\ldots\,]}}

% Lecture douteuse : le mot est donné, et signalé comme incertain.
\newcommand{\uncertain}[1]{\uline{#1}}

% Ajout de l'éditeur — une lettre manquante, un mot sous-entendu.
\newcommand{\add}[1]{\textcolor{blue!50!black}{[#1]}}

% Biffé par Grothendieck lui-même. Ce qu'il a rayé fait partie du document :
% une rature dit souvent où la pensée a changé de direction.
\newcommand{\struck}[1]{\sout{#1}}

% Note du transcripteur, sur l'état matériel du feuillet.
\newcommand{\note}[1]{\par\smallskip{\small\color{gray!60!black}\textit{[#1]}}\par\smallskip}

% Note marginale de Grothendieck, distincte du corps du texte.
\newcommand{\marginal}[1]{\par\smallskip\noindent\hspace{1em}%
  {\small\color{gray!60!black}#1}\par\smallskip}

% --- Titre ------------------------------------------------------------------

\AtBeginDocument{%
  \begin{center}
    {\large\bfseries Fonds Alexandre Grothendieck --- cote n\textsuperscript{o}~\grothFolder}\\[2pt]
    {\itshape\grothFoldertitle}\\[4pt]
    {\small feuillets \grothFirst--\grothLast\ (lot \grothBatch)}\\[2pt]
    {\footnotesize\color{gray!60!black}Transcription établie d'après le fac-similé numérisé
     par l'Université de Montpellier.\\
     Les lectures incertaines sont soulignées, les illisibles notées [\,\ldots\,],
     les ajouts entre crochets.}
  \end{center}
  \vspace{1em}}

\endinput


\folder{54}
\pages{1}{7}
\dating{[à partir de 1971]}
\watermark{Édition de démonstration\\interprétation personnelle de l'œuvre}
\foldertitle{Dualité en cohomologie des V. A [variétés abéliennes] : notes manuscrites (s.d.) — lecture modernisée du dossier entier}

\begin{document}

\section*{Résumé}

\begin{resume}
Une figure, dessinée deux fois dans deux mondes différents, et la question de
savoir si c'est la même.

Une \emph{variété abélienne} est un objet géométrique dont les points
s'additionnent : une courbe elliptique en est le premier exemple, et en
dimension quelconque elles sont l'analogue algébrique du tore. Chacune,
$A$, en possède une seconde, $B$, sa \emph{duale}, construite à partir d'elle
et de même dimension. La dualité n'est pas une symétrie superficielle : la
duale de la duale est la variété de départ, et presque tout ce qu'on sait dire
de $A$ a une contrepartie sur $B$.

Ces sept feuillets posent une question précise sur cette dualité. La
cohomologie est une machine qui associe à un objet géométrique des espaces
vectoriels mesurant, en gros, ses trous. On dispose donc de deux familles
d'espaces, ceux de $A$ et ceux de $B$. Sont-ils duals les uns des autres, au
sens élémentaire de l'algèbre linéaire — celui où une forme sur l'un est un
vecteur de l'autre ? Et si oui, pourquoi ?

L'idée qui arrive est de ne pas comparer les cohomologies directement, mais de
trouver un objet géométrique dont la cohomologie de $B$ soit simplement
l'espace tangent. Cet objet existe : c'est l'\emph{extension universelle} de
$A$ par un groupe vectoriel, une variété $E$ qui contient un espace vectoriel
et dont le quotient par lui est $A$, choisie de sorte que toute autre
extension du même genre s'en déduise. Son espace tangent, calculé en trois
lignes au feuillet 2, est la cohomologie de de Rham de $B$ — et la suite exacte
qui le décrit est la transposée de celle qui décrit la cohomologie de $A$.
Deux objets transposés l'un de l'autre sont duals : la réponse tombe, et elle
tombe en degré 1 seulement, ce qui suffit parce que tous les autres degrés
s'en déduisent par algèbre extérieure.

Le reste du dossier est la tentative de recommencer la même figure en
caractéristique $p$, là où la géométrie cesse d'être docile. Il y a alors une
autre dualité disponible, celle de Cartier, qui apparie les points d'ordre $p$
de $A$ à ceux de $B$ ; et il y a, à l'intérieur de ces points, une partie
infinitésimale attachée à l'espace tangent. Ce que ces feuillets établissent,
c'est que les deux parties infinitésimales, celle de $A$ et celle de $B$, sont
exactement orthogonales l'une à l'autre pour cet appariement, et que les points
d'ordre $p$ de $A$ s'organisent en une extension canonique de la seconde par la
première. La forme est la même qu'au feuillet 2 : une pièce additive venant de
$B$, un objet au milieu, un quotient dual de la pièce correspondante de $B$.
Que ce soit la même figure et non une ressemblance, il écrit vouloir le savoir
— il s'agit, dit-il, de conjuguer ces extensions et dualités avec ce qui donne
la dualité de Cartier — et la page s'arrête sur deux points d'interrogation de
sa main.

Deux feuillets du dossier appartiennent à un tout autre texte : un tapuscrit
sur la platitude, dont il s'est servi comme papier et qu'il a corrigé au
passage. Les corrections y sont mathématiques et non typographiques : il ôte
une hypothèse à un énoncé et en affaiblit d'autant la conclusion, et il promeut
un corollaire au rang de proposition autonome.

\keywords{abelian variety, dual abelian variety, universal vector extension, de Rham cohomology, Hodge filtration, Hodge cohomology, Serre duality, Cartier duality, Weil pairing, Frobenius kernel, Verschiebung, finite group scheme, Dieudonné module, flattening stratification, local criterion of flatness}
\end{resume}

\pagerange{1}{7}
\section*{Le dossier, sa charpente et ses notations}

Sept feuillets, dont un de garde. Le feuillet 1 ne porte que le titre, de sa
main : \emph{Dualité en Cohomologie de V.A.} Les feuillets 2, 4 et 6 portent la
suite ainsi nommée. Les feuillets 3 et 5 sont deux pages d'un tapuscrit sur la
platitude, sans rapport avec la dualité, qu'il a corrigé à la plume ; le
feuillet 7 est un tapuscrit de Bourbaki sur les chambres et les cloisons, au
dos du feuillet 6, qu'il n'a pas touché.\footnote{L'inventaire de Montpellier
avertit que « les pages au verso des documents sont souvent sans rapport avec
les notes mathématiques », et c'est ici le cas deux fois sur trois. On ne lit
pas le feuillet 7 ; on lit les feuillets 3 et 5, parce que sa plume y travaille
et que ce qu'elle y fait est mathématique.}

L'argument de la suite proprement dite tient en quatre stations :

\begin{itemize}
\item le cadre et les deux espaces tangents (feuillet 2, en tête) ;
\item l'extension vectorielle universelle, et la dualité en degré 1
  (feuillet 2, puis le haut du feuillet 4) ;
\item le passage de ce degré à tous les autres, et ce qui reste à faire
  (feuillet 4) ;
\item la même figure en caractéristique $p$ (bas du feuillet 4, puis
  feuillet 6).
\end{itemize}

Les notations sont fixées ici pour tout ce qui suit. $k$ est un corps, $A$ une
variété abélienne sur $k$ de dimension $d$,\footnote{Le feuillet 2 écrit $n$
pour cet entier et le feuillet 4 écrit $d$ ; c'est le même. On retient $d$,
parce que c'est la seule des deux lettres que le manuscrit prenne la peine de
définir (« où $d = \dim A$ »).} et $B = \widehat{A}$ sa duale, de même
dimension. On pose
\[
  \omega_A = H^0(A, \Omega^1_A), \qquad t_A = \operatorname{Lie} A = \omega_A^{\vee},
\]
et de même pour $B$. On note $D(-)$ le dual de Cartier d'un schéma en groupes
fini, $A[p]$ le noyau de la multiplication par $p$,\footnote{Le manuscrit écrit
${}_pA$, l'indice à gauche ; c'était l'usage. On écrit $A[p]$, qui est celui
d'aujourd'hui.} et $H^{*}_{\mathrm{dR}}$ la cohomologie de de Rham, c'est-à-dire
l'hypercohomologie du complexe $\Omega^{\bullet}_A$.\footnote{Le feuillet 2
l'écrit $H^{*}(A) = \mathbb{H}^{*}(A, \Omega^{*}_A)$, le $H$ et le $\Omega$
soulignés deux fois pour marquer qu'il s'agit du complexe et non du seul
faisceau, et glose « Coh.\ de Hodge--DeRham ». La lecture de l'abréviation
« Coh. » est douteuse sur la page.}

\pagerange{2}{2}
\section*{Le cadre : deux espaces tangents (feuillet 2)}

Tout part d'un fait que le manuscrit pose sans le démontrer, et qui est la
raison d'être de la duale : l'espace tangent de $B$ à l'origine est un groupe
de cohomologie de $A$,
\[
  t_B \;=\; H^1(A, \mathcal{O}_A).
\]
Cela tient à ce que $B$ est la composante neutre du schéma de Picard de $A$ :
déformer un fibré en droites au premier ordre, c'est exactement se donner une
classe dans $H^1(A, \mathcal{O}_A)$. La dualité entre $A$ et $B$ est donc, dès
le départ, une affaire de cohomologie et non seulement de géométrie.

En face, l'espace tangent de $A$ est le dual des différentielles invariantes,
et la dualité de Serre le fait remonter en haut du tableau :
\[
  t_A \;=\; H^0(A, \Omega^1_A)^{\vee} \;\simeq\; H^d(A, \Omega^{d-1}_A).
\]
\footnote{C'est la dualité de Serre appliquée à $\Omega^1_A$, où l'on se sert
de ce que le fibré canonique $\Omega^d_A$ est trivial sur une variété
abélienne — ce qui identifie $(\Omega^1_A)^{\vee}$ à $\Omega^{d-1}_A$. Le
manuscrit écrit l'isomorphisme sans un mot de justification ; l'hypothèse de
trivialité, qu'il ne mentionne pas, est ce qui le rend vrai.}

Il vaut la peine de dire tout de suite ce qui rend le calcul de cohomologie
d'une variété abélienne si maniable, car tout le feuillet 4 en dépend. Le
fibré des $1$-formes y est trivial, $\Omega^1_A \cong \mathcal{O}_A \otimes
\omega_A$, donc $\Omega^p_A \cong \mathcal{O}_A \otimes \Lambda^p \omega_A$ et
\[
  H^q(A, \Omega^p_A) \;\cong\; \Lambda^p \omega_A \otimes \Lambda^q t_B ,
\]
puisque $H^q(A, \mathcal{O}_A) = \Lambda^q H^1(A, \mathcal{O}_A) = \Lambda^q
t_B$. Toute la cohomologie de Hodge est donc lue sur les deux espaces
vectoriels $\omega_A$ et $t_B$, de dimension $d$ chacun.

\pagerange{2}{2}
\section*{L'extension universelle, et la dualité en degré 1 (feuillet 2, feuillet 4 en tête)}

Voici l'idée du dossier, et c'est la bonne.

Les extensions de $A$ par le groupe additif sont classées par
$\operatorname{Ext}^1(A, \mathbb{G}_a) \cong H^1(A, \mathcal{O}_A) = t_B$.
Pour un groupe vectoriel $W(V)$ attaché à un espace vectoriel $V$, on a donc
\[
  \operatorname{Ext}^1\!\left(A,\, W(V)\right) \;\cong\; t_B \otimes V
  \;\cong\; \operatorname{Hom}\!\left(t_B^{\vee},\, V\right).
\]
C'est la formule que porte le feuillet 2, insérée en interligne au-dessus d'une
première rédaction biffée. Un tel groupe de classes étant un groupe de
morphismes \emph{issus} d'un objet fixe, il a un élément universel : celui qui
correspond à l'identité, obtenu en prenant $V = t_B^{\vee} = \omega_B$. D'où
l'\emph{extension vectorielle universelle}
\[
  0 \longrightarrow W(\omega_B) \longrightarrow E(A) \longrightarrow A
    \longrightarrow 0 ,
\]
qui est celle par laquelle toute extension de $A$ par un groupe vectoriel se
factorise, d'une seule manière.\footnote{Le manuscrit écrit le terme de gauche
$W(\check{t}_B)$, avec un caron dont il ne dit rien ; on le lit comme le dual
linéaire, $\check{t}_B = t_B^{\vee}$, seule lecture qui rende la suite exacte
correcte et qui s'accorde avec la ligne suivante du feuillet. La convention est
donc nôtre, et non sienne.}

Passons aux espaces tangents. La suite reste exacte, le terme de gauche est son
propre espace tangent, et l'on obtient
\[
  0 \longrightarrow \omega_B \longrightarrow \operatorname{Lie} E(A)
    \longrightarrow t_A \longrightarrow 0 .
\]

Le pas décisif tient en une observation d'identification. Puisque
$\widehat{B} = A$, la formule du cadre appliquée à $B$ donne
$t_A = \operatorname{Lie} \widehat{B} = H^1(B, \mathcal{O}_B)$. La suite
ci-dessus se relit donc entièrement sur $B$ :
\[
  0 \longrightarrow H^0(B, \Omega^1_B) \longrightarrow \operatorname{Lie} E(A)
    \longrightarrow H^1(B, \mathcal{O}_B) \longrightarrow 0 ,
\]
et l'on reconnaît, terme pour terme, la filtration de Hodge de
$H^1_{\mathrm{dR}}(B)$ — que le manuscrit appelle $S_B$.\footnote{$S_A$ et
$S_B$ ne sont définis nulle part dans le dossier ; il les emploie comme des
choses connues. On les lit comme les suites de la filtration de Hodge, $0 \to
H^0(X,\Omega^1_X) \to H^1_{\mathrm{dR}}(X) \to H^1(X,\mathcal{O}_X) \to 0$,
seule lecture sous laquelle les deux affirmations qu'il en fait — que la suite
tangente \emph{est} $S_B$, et que $S_B$ est la transposée de $S_A$ — soient
l'une et l'autre vraies.} D'où sa conclusion, en une ligne :
\[
  \operatorname{Lie} E(A) \;=\; H^1_{\mathrm{dR}}(B).
\]
La cohomologie de de Rham de la duale n'est pas seulement isomorphe à un espace
tangent : elle \emph{est} l'espace tangent d'une variété que $A$ détermine.

Reste la dualité elle-même, et elle est maintenant une transposition. Dualisons
$S_B$ :
\[
  0 \longrightarrow t_A^{\vee} \longrightarrow H^1_{\mathrm{dR}}(B)^{\vee}
    \longrightarrow \omega_B^{\vee} \longrightarrow 0 ,
\]
c'est-à-dire, en réécrivant $t_A^{\vee} = \omega_A$ et $\omega_B^{\vee} = t_B$,
\[
  0 \longrightarrow \omega_A \longrightarrow H^1_{\mathrm{dR}}(B)^{\vee}
    \longrightarrow t_B \longrightarrow 0 ,
\]
qui est $S_A$. Les deux suites sont transposées l'une de l'autre, et par
conséquent
\[
  H^1_{\mathrm{dR}}(A)^{\vee} \;\cong\; H^1_{\mathrm{dR}}(B).
\]
C'est l'encadré du feuillet 2 : « $H^1(A)$ et $H^1(B)$ sont duals l'un de
l'autre ».\footnote{Le lien entre l'extension vectorielle universelle et la
cohomologie de de Rham de la duale, avec la comparaison de la filtration de
Hodge qu'il entraîne, est aujourd'hui associé au mémoire de Mazur et Messing,
\emph{Universal Extensions and One Dimensional Crystalline Cohomology} (1974).
La datation du dossier est celle de l'inventaire, « [à partir de 1971] », et
elle est déduite : rien sur ces feuillets ne les date. On signale la
concordance des sujets et l'on n'en tire aucune conclusion de priorité, que ni
ces pages ni ce fichier ne sont en mesure d'établir.}

\pagerange{4}{4}
\section*{De ce degré à tous les autres (feuillet 4)}

Le feuillet 4 s'ouvre sur la suite de la phrase encadrée, et sur ce qu'il
voudrait de mieux : que $H^{*}(A)$ et $H^{*}(B)$ soient duals en tous degrés,
et non seulement au premier.

Il note aussitôt que ce n'est presque rien de plus, et il a raison. La
cohomologie de Hodge d'une variété abélienne est l'algèbre extérieure sur sa
composante de degré 1 :
\[
  H^{*}(A, \Omega^{*}_A) \;=\; \Lambda^{*}\!\left(\omega_A \oplus t_B\right),
\]
ce qu'il écrit $H^{*}(A,\Omega^{*}_A) = \Lambda\!\left(H^{*}(A,
\Omega^{*}_A)\right)^{(1)}$, l'exposant $(1)$ désignant la partie de degré
$1$.\footnote{Cette identité est immédiate depuis la formule du cadre :
$H^q(A,\Omega^p_A) = \Lambda^p\omega_A \otimes \Lambda^q t_B$ est exactement la
composante de bidegré $(p,q)$ de $\Lambda^{*}(\omega_A \oplus t_B)$. Le
manuscrit dit « est immédiate » et ne la démontre pas ; elle l'est.} Une
dualité entre les composantes de degré 1 se propage donc mécaniquement, et le
calcul donne
\[
  H^q(A, \Omega^p_A)^{\vee} \;\cong\; H^p(B, \Omega^q_B),
\]
puisque le dual de $\Lambda^p\omega_A \otimes \Lambda^q t_B$ est
$\Lambda^p t_A \otimes \Lambda^q \omega_B$, qui est la composante de bidegré
$(q,p)$ de $B$.\footnote{Le feuillet écrit cette dualité entre
$H^p(A,\Omega^q_A)$ et $H^q(B,\Omega^p_A)$. L'échange de $p$ et de $q$ est
juste et il est le point ; l'indice du second $\Omega$, en revanche, est écrit
$A$ là où il faut $B$, et avec $A$ l'énoncé est faux. La lecture de la page est
nette et l'on n'a pas corrigé la transcription : c'est un lapsus de plume dans
une note écrite pour lui-même.} Le même argument, la dégénérescence de la
suite spectrale de Hodge vers de Rham étant acquise pour les variétés
abéliennes, donne $H^i_{\mathrm{dR}}(A)^{\vee} \cong
H^i_{\mathrm{dR}}(B)$ en tous degrés.\footnote{En caractéristique nulle la
dégénérescence est classique ; en caractéristique $p$ elle vaut aussi pour les
variétés abéliennes, résultat que l'on associe aux travaux d'Oda. Le manuscrit
ne la mentionne pas et n'en a pas besoin pour l'énoncé gradué, qui est de
cohomologie de Hodge et non de de Rham ; on la mentionne parce que la lecture
la sollicite pour passer de l'un à l'autre.}

Le feuillet dit alors ce qu'il voudrait à la place, et c'est la seule chose que
le dossier annonce sans jamais la faire : disposer d'une « structure commune »
qui définirait $H^{*}(A)$ et $H^{*}(B)$ d'un seul geste, la dualité n'étant
plus alors qu'une lecture de cette structure. Il ajoute que « cela ne saurait
offrir de difficultés » et passe à autre chose.\footnote{La phrase est reprise
en cours de route sur la page, « il suffirait » biffé, et ses mots de liaison
ne se lisent pas tous ; ce qui précède en donne le sens et non le détail. La
promesse n'est tenue nulle part dans le dossier.}

\pagerange{4}{6}
\section*{La même figure en caractéristique $p$ (feuillets 4 et 6)}

Ici commence la seconde moitié, et elle s'ouvre sur une intention déclarée :
conjuguer les extensions et dualités qui précèdent avec ce qui donne la
\emph{dualité de Cartier}. On se place donc en caractéristique $p > 0$.

\subsection*{L'appariement de Weil, et la partie infinitésimale}

Le point de départ est que la dualité entre $A$ et $B$ se voit sur les points
d'ordre $p$. Le schéma en groupes fini $A[p]$ est de rang $p^{2d}$, et
l'appariement de Weil
\[
  A[p] \times B[p] \longrightarrow \mu_p
\]
est parfait : $B[p] = D(A[p])$, dual de Cartier. C'est ce que le feuillet 4
énonce en disant que $A[p]$ « est dual au sens Cartier » de $B[p]$.

À l'intérieur de $A[p]$ vit une partie purement infinitésimale, le noyau du
Frobenius $F : A \to A^{(p)}$. Elle est de rang $p^d$ et son algèbre de Lie est
$t_A$ ; c'est elle que le manuscrit note $F^A = \operatorname{gr}(t_A)$, le
groupe qu'il attache ainsi à $t_A$.\footnote{Le feuillet 4 écrit
$\operatorname{gr}$ et le
feuillet 6 $\operatorname{Gr}$ ; c'est le même objet et on l'écrit $\ker F_A$.
L'exposant du rang est corrigé sur la page, $p^{2d}$ biffé et remplacé par
$p^{d}$ : la correction est juste, $p^{2d}$ étant le rang de $A[p]$ tout entier
et $p^{d}$ celui de sa partie infinitésimale.}

\subsection*{L'orthogonalité, qui est le résultat du feuillet}

L'affirmation du feuillet 4 est que $\ker F_A$ et $\ker F_B$ sont
\emph{orthogonaux l'un de l'autre} pour l'appariement de Weil. Elle est vraie,
et exactement : chacun est l'annulateur de l'autre, et non seulement contenu
dedans.

La vérification tient en deux lignes une fois posé le dictionnaire de la
dualité de Cartier. Pour un sous-groupe $H \subseteq A[p]$, l'orthogonal
$H^{\perp} \subseteq B[p] = D(A[p])$ est $D(A[p]/H)$. Or l'isogénie $p_A$ se
factorise en $A \xrightarrow{\;F\;} A^{(p)} \xrightarrow{\;V\;} A$, d'où
$A[p]/\ker F_A \cong \ker V_A$, et la dualité échange Frobenius et
Verschiebung : $D(\ker V_A) = \ker F_B$. Donc
\[
  (\ker F_A)^{\perp} \;=\; D\!\left(A[p]/\ker F_A\right) \;=\; D(\ker V_A)
  \;=\; \ker F_B .
\]
Les rangs concordent, $p^d$ de part et d'autre dans un appariement parfait de
rang $p^{2d}$, ce qui était déjà le premier indice.

C'est ce dictionnaire que le feuillet 6 dessine, en portant les deux suites de
Frobenius et de Verschiebung l'une au-dessus de l'autre, « deux suites
transposées l'une de l'autre » :

\begin{tikzcd}[column sep=large]
A \arrow[r, "F"] & A^{(p)} \arrow[r, "V"] & A \\
B & B^{(p)} \arrow[l, "V"] & B \arrow[l, "F"]
\end{tikzcd}

\footnote{Le troisième sommet de la ligne du haut est tracé sur la page comme
un $V$ nu, sans barre ; on le lit $A$ par la symétrie avec la ligne du bas, qui
est complète, et parce que la composée $VF$ est la multiplication par $p$ sur
$A$. La lecture du signe reste douteuse.}

\subsection*{L'extension canonique}

De la factorisation $p = VF$ et du dictionnaire ci-dessus, le feuillet 6 tire
la suite exacte
\[
  0 \longrightarrow \ker F_A \longrightarrow A[p] \longrightarrow
    D(\ker F_B) \longrightarrow 0 ,
\]
qu'il décrit d'un mot : $A[p]$ est l'\emph{extension canonique} de $D(F^B)$ par
$F^A$.

C'est là que la figure se referme, et c'est ce qui fait l'intérêt du dossier.
Comparons avec le feuillet 2 :
\[
  0 \to W(\omega_B) \to E(A) \to A \to 0
  \qquad\text{et}\qquad
  0 \to \ker F_A \to A[p] \to D(\ker F_B) \to 0 .
\]
Même forme : à gauche une pièce additive ou infinitésimale que $B$ fournit, au
milieu l'objet qui porte le « $H^1$ », à droite un quotient dual de la pièce
correspondante de $B$. Dans le premier cas, prendre l'espace tangent donne
$H^1_{\mathrm{dR}}(B)$ et la filtration de Hodge. Dans le second, la théorie de
Dieudonné joue le rôle de l'espace tangent, et le module qu'elle associe à
$A[p]$ se compare de la même façon à la cohomologie de de Rham, la filtration
de Hodge répondant au noyau du Frobenius.\footnote{Cette comparaison est celle
que l'on associe aux travaux d'Oda sur le module de Dieudonné de $A[p]$ et
$H^1_{\mathrm{dR}}(A)$. On l'énonce sans en fixer la variance : selon que l'on
prend le foncteur de Dieudonné covariant ou contravariant, c'est $A$ ou sa
duale qui apparaît, et rien sur ces feuillets ne tranche la convention. Le
dossier ne cite personne et ne mentionne pas la théorie de Dieudonné ; on la
nomme ici parce que c'est le nom sous lequel un lecteur d'aujourd'hui trouvera
la suite.}

Que ce soit là une seule figure et non deux qui se ressemblent, c'est
précisément ce que le feuillet 4 demandait — « conjuguer ces extensions et
dualités avec ce qui donne la dualité de Cartier » — et ce que le dossier ne
démontre pas. Il s'arrête sur un brouillon rapide, une insertion entourée d'un
trait de plume, et deux points d'interrogation de sa main.\footnote{La fin du
feuillet 6 ne se laisse pas lire : les formules portent, les phrases entre
elles non. Le dernier membre lisible est « les correspondances de $A[p]$ ?? »,
et la lecture du mot « correspondances » est douteuse. Les deux points
d'interrogation sont de lui.}

\pagerange{3}{5}
\section*{Un tapuscrit sur la platitude, corrigé au passage (feuillets 3 et 5)}

Les feuillets 3 et 5 n'appartiennent pas à la suite qui précède. Ce sont deux
pages d'un tapuscrit sur la platitude et la liberté des modules, numéroté 4.1
et suivants, écrit à la première personne — « je m'aperçois à l'instant d'une
amélioration dans une autre direction de 4.1., dont je vais reprendre
l'énoncé » — dont le dossier s'est servi comme papier. Sa plume y est passée,
et ce qu'elle y fait n'est pas de la correction typographique.

\subsection*{Une hypothèse ôtée, et la conclusion qu'elle emportait}

La Proposition 4.1 telle que la machine l'a frappée suppose $I$ \emph{nilpotent}
et conclut que $M$ est libre, ou plat, sur $A$. Sa plume biffe « nilpotent » et
récrit la conclusion : pour tout $n$, c'est $M \otimes_A A/I^{n+1}$ qui est
libre, ou plat, sur $A/I^{n+1}$.

L'affaiblissement n'est pas une prudence, c'est ce que la suppression de
l'hypothèse impose exactement. Avec $I$ nilpotent, $I^{n+1} = 0$ pour $n$ assez
grand, donc $A/I^{n+1} = A$ et la conclusion ancienne se retrouve telle quelle :
la nouvelle contient l'ancienne. Sans nilpotence, l'argument ne franchit plus
le passage de $A/I$ à $A$ et ne va pas au-delà des quotients $I$-adiques ; c'est
précisément le pas que la nilpotence servait à faire. On tient donc, après sa
main, l'énoncé suivant :

\begin{quote}
Soient $A$ un anneau, $I$ un idéal de $A$, $(A_i)_{i \in I}$ une famille de
$A$-algèbres telle que l'intersection des noyaux des $A \to A_i$ soit nulle, et
$M$ un $A$-module. Si $M \otimes_A A/I$ est libre sur $A/I$ (resp.\ plat, la
famille d'indices étant alors finie) et si $M \otimes_A A_i$ est libre
(resp.\ plat) sur $A_i$ pour tout $i$, alors pour tout $n$, $M \otimes_A
A/I^{n+1}$ est libre (resp.\ plat) sur $A/I^{n+1}$.
\end{quote}

\footnote{La démonstration de cette forme corrigée n'est pas sur ces deux
feuillets : le texte tapé qui suit l'énoncé traite le cas nilpotent, c'est-à-dire
le cas même que sa plume écarte, et il continue de s'y appuyer une phrase plus
bas. On rapporte donc l'énoncé qu'il laisse sur la page, et la raison pour
laquelle la conclusion devait reculer ; on ne le certifie pas.}

Au bas du feuillet 3 il commence à la main, sous un grand crochet, un
corollaire de cet énoncé corrigé. Il le biffe ligne après ligne et le
brouillon ne se laisse pas reconstituer ; on y relit un anneau, une condition
d'intersection nulle des idéaux, et « $M$ plat sur $A$ ».

\subsection*{Un corollaire promu}

Le feuillet 5 porte ce que l'on appellerait aujourd'hui le sous-schéma de
platification, ou stratification aplatissante : étant donné un morphisme
$f : X \to Y$ de présentation finie et un Module $F$ de présentation finie sur
$X$, il existe, quitte à rétrécir $Y$ à un ouvert non vide, un sous-schéma
fermé $Z \subseteq Y$ tel qu'un changement de base $Y' \to Y$ rende $F'$ plat
sur $Y'$ si et seulement s'il se factorise par $Z$.\footnote{L'énoncé est celui
que l'on trouve, dans sa forme générale, chez Raynaud et Gruson, \emph{Critères
de platitude et de projectivité} (1971) ; sous hypothèses noethériennes et
projectives il était connu plus tôt. Comme au feuillet 2, la datation du
dossier est déduite par l'inventaire et l'on ne conclut rien de la
concordance.}

Trois interventions de sa main sur cet énoncé, et la première est la plus
lourde : le titre « Corollaire 4.1. ter » est biffé et « Proposition A » écrit
au-dessus. L'énoncé cesse d'être un corollaire pour devenir un résultat
autonome. Ce n'est pas une décision prise à cet instant : le feuillet 3 se sert
déjà du nom « Prop.\ A », mais là c'est la machine qui l'a frappé, par-dessus
une référence à 4.1 biffée au ruban. Le tapuscrit se citait donc sous un nom
que son propre titre ne portait pas, et la plume ne fait que fermer l'écart.

Les deux autres corrections sont des hypothèses. « $Y$ affine » tombe au profit
de « $Y \neq \emptyset$ » — l'énoncé n'a pas besoin d'affinité, seulement de
quelque chose à rétrécir. Et $Z$ reçoit d'être \emph{de présentation finie sur
$Y$}, sans quoi le sous-schéma obtenu ne serait pas un objet avec lequel on
peut travailler dans la catégorie où l'énoncé se place.

\end{document}
